\documentclass[aspectratio=169]{beamer}

\usepackage{listings}
\usepackage{xcolor}

\definecolor{codegray}{rgb}{0.5,0.5,0.5}
\definecolor{codepurple}{rgb}{0.58,0,0.82}
\definecolor{backcolour}{rgb}{0.95,0.95,0.92}

\lstdefinestyle{customc}{
  backgroundcolor=\color{backcolour},
  commentstyle=\color{codegray}\itshape,
  keywordstyle=\color{blue}\bfseries,
  numberstyle=\tiny\color{gray},
  stringstyle=\color{codepurple},
  basicstyle=\ttfamily\footnotesize,
  breakatwhitespace=false,
  breaklines=true,
  captionpos=b,
  keepspaces=true,
  numbers=left,
  numbersep=5pt,
  showspaces=false,
  showstringspaces=false,
  showtabs=false,
  tabsize=2,
  frame=single,
  language=C++
}

\lstset{style=customc}

\usepackage[spanish]{babel}
\usepackage[utf8]{inputenc}
\usepackage[T1]{fontenc}
\usepackage{lmodern}
\usepackage{amsmath, amsfonts}
\usepackage{hyperref}

\usetheme{Madrid}

\title{Aritmética modular}
\author{Sebastian Mestre}
\date{}

\begin{document}

\begin{frame}
  \titlepage
\end{frame}

\section{Aritmética modular}

\begin{frame}{¿Qué es la aritmética modular?}
  \begin{itemize}
    \item Sistema de numeración construido sobre el \textbf{resto} de la división.
    \item Dos números son \textbf{congruentes módulo \(m\)} si su diferencia es un múltiplo de \(m\).
  \end{itemize}
  \begin{block}{Definición}
    Dados enteros \(a,b\) y \(m \ge 1\),
    \[
      a \equiv b \pmod m
      \quad \Leftrightarrow \quad
      a - b = k \cdot m \text{ para algún } k \in \mathbb{Z}.
    \]
  \end{block}
  \begin{itemize}
    \item En particular, sumar \(m\) no cambia la clase:
      \[
        a + m \equiv a \pmod m.
      \]
  \end{itemize}
\end{frame}

\begin{frame}{Propiedades de la congruencia}
  \begin{itemize}
    \item La relación \(\equiv \pmod m\) es una \textbf{relación de equivalencia}:
      \begin{itemize}
        \item Reflexiva: \(a \equiv a \pmod m\).
        \item Simétrica: si \(a \equiv b \pmod m\), entonces \(b \equiv a \pmod m\).
        \item Transitiva: si \(a \equiv b \pmod m\) y \(b \equiv c \pmod m\), entonces \(a \equiv c \pmod m\).
      \end{itemize}
    \item Se comporta bien con las operaciones básicas:
      \[
        \begin{aligned}
          a &\equiv b \pmod m \\
          c &\equiv d \pmod m
        \end{aligned}
        \quad \Rightarrow \quad
        \begin{cases}
          a + c \equiv b + d \pmod m,\\
          a - c \equiv b - d \pmod m,\\
          a \cdot c \equiv b \cdot d \pmod m.
        \end{cases}
      \]
  \end{itemize}
\end{frame}

\begin{frame}{Clases de equivalencia}
  \begin{itemize}
    \item Para cada entero \(a\), consideramos el conjunto de todos los números congruentes con \(a\) módulo \(m\).
    \item Este conjunto se llama \textbf{clase de equivalencia} de \(a\) módulo \(m\):
      \[
        [a] = \{\, b \in \mathbb{Z} \mid b \equiv a \pmod m \,\}.
      \]
  \end{itemize}
  \begin{block}{Ejemplo: \(m = 5\)}
    \[
    \begin{aligned}
      [0] &= \{\dots, -5, 0, 5, 10, 15, \dots\}\\
      [1] &= \{\dots, -4, 1, 6, 11, 16, \dots\}\\
      [2] &= \{\dots, -3, 2, 7, 12, 17, \dots\}\\
      [3] &= \{\dots, -2, 3, 8, 13, 18, \dots\}\\
      [4] &= \{\dots, -1, 4, 9, 14, 19, \dots\}
    \end{aligned}
    \]
    Notar que \([5] = [0]\), es la misma clase.
  \end{block}
\end{frame}

\begin{frame}{El conjunto \( \mathbb{Z}_m \)}
  \begin{itemize}
    \item Hay exactamente \(m\) clases de equivalencia módulo \(m\):
      \[
        [0], [1], \dots, [m-1].
      \]
    \item Al conjunto de estas clases se lo denota típicamente por \(\mathbb{Z}_m\):
      \[
        \mathbb{Z}_m = \{ [0], [1], \dots, [m-1] \}.
      \]
    \item Como la suma, la resta y el producto \textbf{respetan} la congruencia, podemos definir operaciones entre clases:
      \[
        [x] + [y] = [x + y],\quad
        [x] - [y] = [x - y],\quad
        [x] \cdot [y] = [x \cdot y].
      \]
  \end{itemize}
\end{frame}

\section{Representación canónica}

\begin{frame}{Representación canónica}
  \begin{itemize}
    \item Para todo entero \(x\), existe un único entero \(r \in [0, m-1]\) tal que
      \[
        x \equiv r \pmod m.
      \]
    \item A este \(r\) lo llamamos \textbf{representación canónica} de \(x\) módulo \(m\).
    \item Si \(x\) es positivo, \(r\) es simplemente el resto de dividir \(x\) por \(m\).
    \item Es muy cómodo trabajar siempre con los números en su representación canónica.
  \end{itemize}
\end{frame}

\begin{frame}[fragile]{Cuidado con \texttt{\%} en C++}
  \begin{itemize}
    \item En C++, cuando \(x\) es negativo, \texttt{x \% m} \textbf{no} da la representación canónica (puede ser negativa).
    \item Podemos corregirlo de forma estándar:
  \end{itemize}
  \begin{block}{Función de representación canónica}
  \begin{lstlisting}[style=customc]
int const m = 1000000007; // muy importante que sea constante

int canonica(int x) {
  return (x % m + m) % m;
}
  \end{lstlisting}
  \end{block}
\end{frame}

\section{Operaciones básicas}

\begin{frame}[fragile]{Operaciones en representación canónica}
  \begin{itemize}
    \item Si mantenemos todos los números en \([0, m-1]\), las operaciones básicas se pueden implementar de forma segura.
    \item Truco: una versión más rápida que \texttt{canonica} cuando sabemos que el resultado está en \([0, 2m-1]\).
  \end{itemize}
  \begin{block}{Implementación}
  \begin{lstlisting}[style=customc]
// igual que canonica en el intervalo [0, 2m-1], y más rápido
int nm(int x) {
  return x - m * (x >= m);
}

// si x e y están en [0, m-1], x + y está en [0, 2m-1]
int add(int x, int y) {
  return nm(x + y);
}

// si x e y están en [0, m-1], x + m - y está en [0, 2m-1]
int sub(int x, int y) {
  return nm(x + m - y);
}

// cuidado: x * y puede desbordar int, usamos long long
int mul(int x, int y) {
  return (long long)x * y % m;
}
  \end{lstlisting}
  \end{block}
\end{frame}

\section{Otras operaciones}

\begin{frame}[fragile]{Potenciación modular}
  \begin{itemize}
    \item Necesitamos calcular \(x^y \bmod m\) muchas veces en programación competitiva.
    \item Usamos \textbf{exponenciación binaria} para obtener costo \(O(\log y)\).
  \end{itemize}
  \begin{block}{Implementación}
  \begin{lstlisting}[style=customc]
// O(log y)
int pot(int x, int y) {
  if (y == 0) return 1;
  if (y % 2 == 0) return pot(mul(x, x), y / 2);
  return mul(x, pot(x, y - 1));
}
  \end{lstlisting}
  \end{block}
\end{frame}

\begin{frame}[fragile]{Inverso multiplicativo}
  \begin{itemize}
    \item Si \(m\) es primo, podemos usar el \textbf{pequeño teorema de Fermat}:
      \[
        x^{m-1} \equiv 1 \pmod m \quad \Rightarrow \quad
        x^{m-2} \equiv x^{-1} \pmod m.
      \]
    \item Entonces, el inverso multiplicativo se puede calcular en \(O(\log m)\).
  \end{itemize}
  \begin{block}{Implementación}
  \begin{lstlisting}[style=customc]
int inv(int x) {
  return pot(x, m - 2);
}
  \end{lstlisting}
  \end{block}
\end{frame}

\begin{frame}[fragile]{Trucos con inversos y factoriales}
  \begin{itemize}
    \item A veces queremos todos los inversos \(1, 2, \dots, n\) en \(O(n)\).
    \item También es común precomputar factoriales y sus inversos para combinatoria.
  \end{itemize}
  \begin{block}{Implementaciones típicas}
  \begin{lstlisting}[style=customc]
int inv[n + 1];
inv[1] = 1;
for (int i = 2; i <= n; i++) {
  inv[i] = mul(inv[m % i], m - m / i);
}

int fac[n + 1], ifac[n + 1];
fac[0] = 1;
forn(i, n) fac[i + 1] = mul(fac[i], i + 1);  // factoriales
ifac[n] = inv(fac[n]);
dforn(i, n) ifac[i] = mul(ifac[i + 1], i + 1); // inversos factoriales
forr(i, 1, n) inv[i] = mul(ifac[i], fac[i - 1]); // inversos 1..n
  \end{lstlisting}
  \end{block}
\end{frame}

\section*{Cierre}

\begin{frame}{Ideas clave}
  \begin{itemize}
    \item La aritmética modular trabaja con clases de equivalencia definidas por el resto módulo \(m\).
    \item Mantener siempre la \textbf{representación canónica} simplifica las implementaciones y evita errores.
    \item Con operaciones básicas (\texttt{add}, \texttt{sub}, \texttt{mul}, \texttt{pot}, \texttt{inv}) podemos resolver muchos problemas de combinatoria y conteo.
  \end{itemize}
\end{frame}

\begin{frame}{Preguntas}
  \centering
  \Huge
  \textbf{Q\&A}
\end{frame}

\end{document}


